\documentclass[11pt]{article}

\usepackage{amsmath,amssymb}
\usepackage{geometry}
\geometry{margin=1in}

\title{A Geometric Intersection View of Eigenvalue Problems and Implicit Surface Reconstruction (MCA\_WLT)}
\author{Chaday Tai}
\date{}

\begin{document}

\maketitle

\begin{abstract}
We present a simple geometric interpretation linking eigenvalue problems in linear algebra and implicit surface reconstruction in computational geometry.

We show that eigenvectors, gradients, and surface intersection points can all be interpreted as geometric intersection structures in $\mathbb{R}^3$.

We also introduce MCA\_WLT (Marching Cubes Algorithm Without Lookup Table), a method that reconstructs implicit surfaces using direct edge-based intersection instead of lookup tables.

The main contribution is a unified geometric viewpoint rather than a new numerical optimization method.
\end{abstract}

\section{Eigenvalue Problem as Geometry}

The eigenvalue equation is:

\[
Av = \lambda v
\]

Rewriting:

\[
(A - \lambda I)v = 0
\]

Let $v = (x,y,z)$. This becomes three plane equations:

\[
(A - \lambda I)v = 0
\]

Each row represents a plane passing through the origin.

Thus, eigenvectors correspond to the intersection of planes in $\mathbb{R}^3$.

\section{Geometric Meaning}

The eigenvector is a direction that satisfies all constraints:

\[
(A - \lambda I)v = 0
\]

It represents the intersection direction of multiple geometric planes.

\section{Implicit Surface and Gradient}

For an implicit surface:

\[
f(x,y,z)=0
\]

The gradient:

\[
\nabla f
\]

is perpendicular to the surface and defines the surface normal direction.

\section{MCA\_WLT Method}

Given a scalar field $f(x,y,z)$ on a voxel grid:

A surface is detected on cube edges where:

\[
f(p_1)\cdot f(p_2) < 0
\]

The intersection point is computed by linear interpolation:

\[
p = p_1 + t(p_2 - p_1), \quad
t = \frac{f(p_1)}{f(p_1)-f(p_2)}
\]

This avoids lookup tables used in classical Marching Cubes.

\section{Unified Interpretation}

Eigenvectors, gradients, and MCA\_WLT intersections all share the same idea:

\begin{quote}
They are geometric intersection problems in $\mathbb{R}^3$.
\end{quote}

\section{Conclusion}

We provide a simple geometric framework that connects linear algebra and surface reconstruction through intersection-based interpretation.

This suggests that many problems in linear algebra and geometry can be unified under a common geometric viewpoint.

\end{document}